\section{Distância em \texorpdfstring{$\mathbb R^n$}{Rn}}
A distância euclidiana usual entre dois pontos $p,q\in\mathbb R^n$ é motivada pelo Teorema de Pitágoras:
ela mede o comprimento do vetor deslocamento entre esses dois pontos.

\begin{definition}
    Sejam $p=(p_1,\ldots,p_n)$ e $q=(q_1,\ldots,q_n)$ elementos de $\mathbb R^n$.
    A \emph{distância euclidiana usual}\index{distância euclidiana}\index{métrica euclidiana|see{distância euclidiana}} entre $p$ e $q$ é definida como:
    \begin{equation*}
        d(p, q) = \|p-q\| = \sqrt{(p-q) \cdot (p-q)}=\sqrt{\sum_{i=1}^n (p_i - q_i)^2}.
    \end{equation*}
\end{definition}

Assim, medir a distância entre $p$ e $q$ é medir o comprimento do vetor $p-q$.
Como esse vetor representa um deslocamento entre os dois pontos, a definição acima estende a ideia pitagórica de distância.

No caso $n=1$, identificando $\mathbb R^1$ com $\mathbb R$, se $p,q\in\mathbb R$, então
\[
    d(p,q)=|p-q|.
\]

No caso $n=2$, se $p=(x_1,y_1)$ e $q=(x_2,y_2)$, então
\[
    d(p,q)=\sqrt{(x_1-x_2)^2+(y_1-y_2)^2},
\]
que é exatamente a fórmula da distância no plano cartesiano.

Já no caso $n=3$, se $p=(x_1,y_1,z_1)$ e $q=(x_2,y_2,z_2)$, recuperamos, analogamente, a fórmula usual da distância no espaço tridimensional:
\[
    d(p,q)=\sqrt{(x_1-x_2)^2+(y_1-y_2)^2+(z_1-z_2)^2}.
\]

Para $n\geq 4$, a mesma fórmula será tomada como a extensão natural da distância pitagórica para pontos com $n$ coordenadas.

\begin{example}\label{ex:distanciaRn}
    Se $p=(1,2)$ e $q=(4,6)$, então
    \[
        d(p,q)=\sqrt{(1-4)^2+(2-6)^2}
        =\sqrt{9+16}=5.
    \]
    Geometricamente, estamos calculando o comprimento do vetor
    \[
        p-q=(-3,-4),
    \]
    que representa um deslocamento de $q$ até $p$.
\end{example}

\subimport{distanciaRn}{figura}

As propriedades abaixo formalizam fatos esperados de qualquer boa noção de distância:
a distância nunca é negativa; dois pontos têm distância zero exatamente quando são o mesmo ponto;
a ordem dos pontos não altera a distância entre eles;
e ir diretamente de $p$ até $r$ nunca é mais longo do que ir de $p$ até $q$ e depois de $q$ até $r$.
\begin{proposition}\label{prop:distanciaRn}
    Sejam $p, q, r \in \mathbb R^n$. Então:
    \begin{enumerate}
        \item $d(p, q)\geq 0$, e $d(p, q) = 0$ se, e somente se, $p=q$.
        \item $d(p, q) = d(q, p)$ (simetria);
        \item $d(p, r) \leq d(p, q) + d(q, r)$ (desigualdade triangular).
    \end{enumerate}
\end{proposition}
\begin{proof}
    Vamos verificar cada uma das propriedades.
    \begin{enumerate}
        \item Como $d(p,q)=\|p-q\|$, temos $d(p,q)\geq 0$. Além disso, $d(p, q) = 0$ se, e somente se, $\|p-q\|=0$, o que ocorre se, e somente se, $p-q=0$, ou seja, $p=q$.
        \item Como $q-p=-(p-q)$, temos
        \[
            d(q,p)=\|q-p\|=\|-(p-q)\|=|-1|\|p-q\|=\|p-q\|=d(p,q).
        \]
        \item Como $p-r=(p-q)+(q-r)$, temos:
        \begin{align*}
            d(p, r) &= \|p-r\| = \|(p-q)+(q-r)\| \\
            &\leq \|p-q\| + \|q-r\| = d(p, q) + d(q, r).
        \end{align*}
    \end{enumerate}
\end{proof}

\begin{exercise}
    Calcule as seguintes distâncias:
    \begin{enumerate}[label=(\alph*)]
        \item $d(-3,5)$ em $\mathbb R$;
        \item $d((1,2),(4,6))$ em $\mathbb R^2$;
        \item $d((0,0),(3,4))$ em $\mathbb R^2$;
        \item $d((1,0,-1,2),(3,-1,2,2))$ em $\mathbb R^4$.
    \end{enumerate}
\end{exercise}

\begin{exercise}
    Sejam $p,q,a\in\mathbb R^n$.
    Mostre que
    \[
        d(p+a,q+a)=d(p,q).
    \]
    Interprete geometricamente essa igualdade.
\end{exercise}

\begin{exercise}
    Sejam $p,q\in\mathbb R^n$ e $\alpha\in\mathbb R$.
    Mostre que
    \[
        d(\alpha p,\alpha q)=|\alpha|d(p,q).
    \]
\end{exercise}

\begin{exercise}
    Sejam $p,q,r\in\mathbb R^n$.
    Use a desigualdade triangular para mostrar que
    \[
        d(p,r)-d(q,r)\leq d(p,q).
    \]
    Em seguida, conclua que
    \[
        |d(p,r)-d(q,r)|\leq d(p,q).
    \]
\end{exercise}

\begin{exercise}
    Determine o conjunto de todos os pontos $x\in\mathbb R$ tais que $d(x,3)<2$.
\end{exercise}
